\clearpage
\section*{Problem 7}
\addcontentsline{toc}{section}{Problem 7}
\markboth{Problem 7}{Problem 7}

\textbf{Problem Statement (Textbook 4.2.5):} 

Prove that an upper or lower triangular matrix is nonsingular if and only if its diagonal elements are all different from $0$.

\noindent\rule{\textwidth}{0.4pt}

\vspace{1em}

\textbf{Solution:}

We prove the result for upper triangular matrices; the lower triangular case is analogous.

\textbf{($\Rightarrow$) Nonsingular implies nonzero diagonal:}

Suppose $U$ is upper triangular and nonsingular. For a triangular matrix, $\det(U) = \prod_{i=1}^{n} u_{ii}$ (the product of diagonal elements). Since $U$ is nonsingular, $\det(U) \neq 0$, which implies each diagonal element $u_{ii} \neq 0$. $\square$

\vspace{1.5em}

\textbf{($\Leftarrow$) Nonzero diagonal implies nonsingular:}

Suppose $U$ is upper triangular with all diagonal elements nonzero. We show $U\mathbf{x} = \mathbf{0}$ has only the trivial solution $\mathbf{x} = \mathbf{0}$ (thus $U$ is nonsingular).

Using back substitution on $U\mathbf{x} = \mathbf{0}$:
\begin{itemize}
    \item Row $n$: $u_{nn}x_n = 0$. Since $u_{nn} \neq 0$, we get $x_n = 0$.
    \item Row $n-1$: $u_{n-1,n-1}x_{n-1} + u_{n-1,n}x_n = 0$. Since $x_n = 0$ and $u_{n-1,n-1} \neq 0$, we get $x_{n-1} = 0$.
    \item Continuing upward, each row $i$ gives: $u_{ii}x_i + \sum_{j=i+1}^{n} u_{ij}x_j = 0$. Since all $x_j = 0$ for $j > i$ and $u_{ii} \neq 0$, we get $x_i = 0$.
\end{itemize}

By induction, $\mathbf{x} = \mathbf{0}$ is the only solution, so $U$ is nonsingular. $\square$

\vspace{1em}

\textbf{Example:} Consider the upper triangular matrix:
$$U = \begin{bmatrix} 2 & 3 & 1 \\ 0 & -1 & 4 \\ 0 & 0 & 5 \end{bmatrix}$$

All diagonal elements $(2, -1, 5)$ are nonzero, so $U$ is nonsingular with $\det(U) = 2 \cdot (-1) \cdot 5 = -10 \neq 0$.

Conversely, consider:
$$V = \begin{bmatrix} 2 & 3 & 1 \\ 0 & 0 & 4 \\ 0 & 0 & 5 \end{bmatrix}$$

The diagonal has a zero at position $(2,2)$, so $\det(V) = 2 \cdot 0 \cdot 5 = 0$ and $V$ is singular. Indeed, $V\mathbf{x} = \mathbf{0}$ has nontrivial solutions such as $\mathbf{x} = [3, -2, 0]^T$.

\vspace{2em}
\noindent\rule{\textwidth}{0.4pt}
